\documentclass[12pt]{article}

\input{~/School/Notes/header.tex}

\title{Домашнее задание по алгебре (к 13.10)}
\author{Денисов Илья}
\date{}

\begin{document}
	\maketitle

	%==================================================
	\section*{Задача 1}

	\noindent
	\textbf{Условие.} Группа $G$ действует на множестве $M$, $|G|=55$,
	$|M|=39$. Докажите, что найдётся $m\in M$ такой, что $gm=m$ для любого
	$g\in G$.

	\begin{proof}
		Рассмотрим орбиты $Gm=\{gm\mid g\in G\}$. Орбиты либо совпадают,
		либо не пересекаются, поэтому $M$ разбивается в объединение
		непересекающихся орбит, и
		\[
			39=|M|=\sum_{\text{орбиты}}|Gm|.
		\]

		Найдём возможные размеры орбит. Зафиксируем $m\in M$ и обозначим
		$\operatorname{Stab}(m)=\{g\in G\mid gm=m\}$~--- подгруппу в $G$.
		Для $g,g'\in G$ имеем
		\[
			gm=g'm
			\iff g^{-1}g'm=m
			\iff g^{-1}g'\in\operatorname{Stab}(m)
			\iff g\operatorname{Stab}(m)=g'\operatorname{Stab}(m).
		\]
		Значит, отображение $g\operatorname{Stab}(m)\mapsto gm$ корректно задаёт
		биекцию между левыми смежными классами по $\operatorname{Stab}(m)$ и
		орбитой $Gm$. По теореме Лагранжа
		\[
			|Gm|=[G:\operatorname{Stab}(m)]=\frac{|G|}{|\operatorname{Stab}(m)|}.
		\]
		Итак, размер каждой орбиты делит $|G|=55=5\cdot11$, то есть равен
		$1$, $5$, $11$ или $55$. Орбита не может быть больше $M$, а
		$55>39$, поэтому
		\[
			|Gm|\in\{1,5,11\}.
		\]

		Допустим, что орбит размера $1$ нет. Тогда все орбиты имеют размер $5$
		или $11$, и для некоторых целых $a,b\geqslant0$
		\[
			39=5a+11b.
		\]
		Так как $11b\leqslant39$, возможны лишь $b=0,1,2,3$, и тогда
		$39-11b$ равно соответственно $39,\ 28,\ 17,\ 6$. Ни одно из этих
		чисел не делится на $5$, то есть $a$ целым быть не может. Противоречие.

		Следовательно, есть орбита размера $1$: $Gm=\{m\}$. Это в точности
		означает, что $gm=m$ для любого $g\in G$.
	\end{proof}

	%==================================================
	\section*{Задача 2}

	\noindent
	\textbf{Условие.} $G$~--- группа, $H$~--- её подгруппа. Докажите, что
	можно придумать такое действие $G$, что стабилизатор одной из точек
	равен $H$.

	\begin{proof}
		Пусть $M=G/H=\{xH\mid x\in G\}$~--- множество левых смежных классов
		по $H$. Зададим действие $G$ на $M$ формулой
		\[
			g\cdot(xH)=(gx)H.
		\]

		\emph{Корректность.} Формула не зависит от выбора представителя
		класса: если $xH=yH$, то $y=xh$ для некоторого $h\in H$, и
		$gyH=gxhH=gxH$. Иными словами, $g\cdot C=\{gc\mid c\in C\}$~---
		образ множества $C$ при умножении слева на $g$, и он снова является
		смежным классом.

		\emph{Аксиомы действия.}
		\begin{enumerate}
			\item $e\cdot(xH)=(ex)H=xH$.
			\item $(g_1g_2)\cdot(xH)=(g_1g_2x)H=g_1\cdot\bigl((g_2x)H\bigr)
				=g_1\cdot\bigl(g_2\cdot(xH)\bigr)$.
		\end{enumerate}

		\emph{Стабилизатор точки $H=eH\in M$.} По определению
		\[
			\operatorname{Stab}(H)=\{g\in G\mid gH=H\}.
		\]
		Если $g\in H$, то $gH=H$, так как $H$~--- подгруппа. Обратно, если
		$gH=H$, то $g=ge\in gH=H$. Значит, $\operatorname{Stab}(H)=H$.
	\end{proof}

	\noindent
	\textbf{Замечание.} Стабилизатор произвольной точки $xH$ равен
	$xHx^{-1}$: действительно, $gxH=xH\iff x^{-1}gx\in H$. Поэтому
	именно точка $eH$ даёт ровно $H$.

	%==================================================
	\section*{Задача 3}

	\noindent
	\textbf{Условие.} Некоторая группа $G$ перестановок на $n$ элементах
	(то есть подгруппа $S_n$) обладает следующим свойством: для любых
	$a\ne b$, $c\ne d$ (любые другие пары букв могут совпадать) существует
	единственная $\pi\in G$ такая, что $\pi(a)=c$, $\pi(b)=d$. Докажите, что
	в $G$ есть подгруппа а)~порядка $n-1$; б)~порядка $n$.

	\medskip
	\noindent
	Будем считать $n\geqslant2$ (при $n=1$ подгруппы порядка $n-1=0$ не
	бывает). Произведение перестановок~--- композиция:
	$(\pi\sigma)(x)=\pi(\sigma(x))$. Заданное свойство будем называть
	$(\ast)$.

	\begin{proof}[Предварительное наблюдение]
		\emph{Любая перестановка $\pi\in G$, $\pi\ne\mathrm{id}$, имеет не более
		одной неподвижной точки.}

		Действительно, пусть $\pi(a)=a$ и $\pi(b)=b$ при $a\ne b$. Тождественная
		перестановка тоже удовлетворяет условию $\mathrm{id}(a)=a$,
		$\mathrm{id}(b)=b$. По единственности в $(\ast)$ получаем
		$\pi=\mathrm{id}$.
	\end{proof}

	\subsection*{а) Подгруппа порядка $n-1$}

	\begin{proof}
		Рассмотрим стабилизатор точки $1$:
		\[
			H=\operatorname{Stab}(1)=\{\pi\in G\mid\pi(1)=1\}.
		\]
		Это подгруппа в $G$. Найдём её порядок.

		Для каждого $d\in\{2,\ldots,n\}$ в силу $(\ast)$ (при $a=1$, $b=2$,
		$c=1$) существует единственная $\pi_d\in G$ с
		$\pi_d(1)=1$, $\pi_d(2)=d$. Все $\pi_d$ лежат в $H$ и попарно различны
		(у них разные значения в точке $2$). С другой стороны, любая
		$\pi\in H$ удовлетворяет $\pi(2)\ne\pi(1)=1$, поэтому
		$\pi(2)=d$ для некоторого $d\in\{2,\ldots,n\}$, и по единственности
		$\pi=\pi_d$. Следовательно,
		\[
			|H|=n-1.
		\]
	\end{proof}

	\subsection*{б) Подгруппа порядка $n$}

	\noindent
	Назовём перестановку \emph{свободной}, если у неё нет неподвижных точек.
	Положим
	\[
		N=\{\mathrm{id}\}\cup\{\pi\in G\mid\pi\text{ свободна}\}.
	\]
	Докажем, что $N$~--- подгруппа порядка $n$.

	\begin{proof}[Шаг 1: для любых $a,c$ в $N$ есть ровно одна перестановка с $\rho(a)=c$]
		\emph{Случай $a=c$.} Перестановка $\mathrm{id}$ подходит, а остальные
		элементы $N$ свободны, то есть не фиксируют $a$. Значит, подходит
		ровно одна.

		\emph{Случай $a\ne c$.} Тождественная перестановка не подходит, поэтому
		нужно найти число свободных $\pi\in G$ с $\pi(a)=c$.

		Сначала посчитаем все $\pi\in G$ с $\pi(a)=c$. Выберем $b\ne a$. Для
		каждого $d\ne c$ по $(\ast)$ есть единственная $\pi_d\in G$ с
		$\pi_d(a)=c$, $\pi_d(b)=d$. Любая $\pi$ с $\pi(a)=c$ имеет
		$\pi(b)\ne c$ (иначе $\pi(a)=\pi(b)$), поэтому равна некоторой
		$\pi_d$. Разным $d$ отвечают разные $\pi_d$. Итак, всего таких
		перестановок $n-1$.

		Теперь посчитаем среди них те, у которых есть неподвижная точка $x$.
		Здесь $x\ne a$ (так как $\pi(a)=c\ne a$) и $x\ne c$ (иначе
		$\pi(c)=c=\pi(a)$ при $c\ne a$). Для каждого из $n-2$ чисел
		$x\notin\{a,c\}$ по $(\ast)$ существует единственная $\sigma_x\in G$ с
		$\sigma_x(a)=c$, $\sigma_x(x)=x$; других перестановок с неподвижной
		точкой нет. Все $\sigma_x$ различны: $\sigma_x\ne\mathrm{id}$, поэтому
		по предварительному наблюдению $x$~--- её единственная неподвижная
		точка. Значит, перестановок с $\pi(a)=c$ и с неподвижной точкой ровно
		$n-2$, а свободных
		\[
			(n-1)-(n-2)=1.
		\]
	\end{proof}

	\begin{proof}[Шаг 2: $|N|=n$]
		По шагу 1 отображение $\rho\mapsto\rho(1)$ из $N$ в $\{1,\ldots,n\}$
		инъективно (единственность) и сюръективно (существование). Значит,
		$|N|=n$.
	\end{proof}

	\begin{proof}[Шаг 3: $N$~--- подгруппа]
		Ясно, что $\mathrm{id}\in N$. Если $\rho\in N$, то $\rho^{-1}\in N$:
		равенство $\rho^{-1}(x)=x$ равносильно $\rho(x)=x$, поэтому
		$\rho^{-1}$ свободна вместе с $\rho$.

		Пусть $\pi,\sigma\in N$. Покажем, что $\pi\sigma\in N$. Если одна из
		перестановок тождественна, то это очевидно, поэтому пусть обе свободны.
		Если $\pi\sigma$ свободна, то $\pi\sigma\in N$ по определению. Иначе
		существует $x$ с $\pi(\sigma(x))=x$. Положим $y=\sigma(x)$. Тогда
		$\pi(y)=x$, то есть $\pi^{-1}(x)=y$. Итак, обе перестановки
		$\sigma$ и $\pi^{-1}$ лежат в $N$ и переводят $x$ в $y$. По шагу 1
		они совпадают: $\sigma=\pi^{-1}$. Тогда $\pi\sigma=\mathrm{id}\in N$.
	\end{proof}

	\noindent
	Итак, $N$~--- подгруппа порядка $n$.

	%==================================================
	\section*{Задача 4}

	\noindent
	\textbf{Условие.} Пусть $a$~--- корень многочлена $x^3+px^2+qx+r$.
	Напишите многочлен, корнем которого является а)~$\dfrac1{a+1}$;
	б)~$a^2+a+1$ (в ответе, разумеется, не должно явным образом
	присутствовать $a$).

	\medskip
	\noindent
	Итак, $a^3+pa^2+qa+r=0$, коэффициенты $p,q,r$ заданы.

	\subsection*{а) Корень $b=\dfrac1{a+1}$}

	\begin{proof}[Решение]
		Число $b$ определено только при $a\ne-1$; будем считать это
		выполненным. Из $b(a+1)=1$ получаем $b\ne0$ и
		\[
			a=\frac{1-b}{b}.
		\]
		Подставим в $a^3+pa^2+qa+r=0$ и домножим равенство на $b^3$:
		\[
			(1-b)^3+p\,b(1-b)^2+q\,b^2(1-b)+r\,b^3=0.
		\]
		Раскроем скобки:
		\begin{align*}
			(1-b)^3&=1-3b+3b^2-b^3,\\
			p\,b(1-b)^2&=pb-2pb^2+pb^3,\\
			q\,b^2(1-b)&=qb^2-qb^3.
		\end{align*}
		Соберём коэффициенты при степенях $b$:
		\[
			(p-q+r-1)\,b^3+(q-2p+3)\,b^2+(p-3)\,b+1=0.
		\]
		Значит, $b$~--- корень многочлена
		\[
			\boxed{g(y)=(p-q+r-1)\,y^3+(q-2p+3)\,y^2+(p-3)\,y+1}
		\]
	\end{proof}

	\noindent
	\textbf{Замечания.} Свободный член $g$ равен $1$, поэтому $g\ne0$ при любых
	$p,q,r$. Старший коэффициент $p-q+r-1=f(-1)$, где
	$f(x)=x^3+px^2+qx+r$. Если $f(-1)\ne0$, то $g$ можно поделить на него и
	получить унитарный многочлен.

	\noindent
	\emph{Проверка.} Для $f=x^3-1$ ($p=q=0$, $r=-1$) и $a=1$ имеем
	$b=\tfrac12$, а $g(y)=-2y^3+3y^2-3y+1$, и
	$g(\tfrac12)=-\tfrac14+\tfrac34-\tfrac32+1=0$.

	\subsection*{б) Корень $c=a^2+a+1$}

	\begin{proof}[Решение]
		Любая степень $a$ выражается через $1,a,a^2$ с помощью
		$a^3=-pa^2-qa-r$. В частности,
		\[
			a^4=a\cdot a^3=-pa^3-qa^2-ra
			=(p^2-q)a^2+(pq-r)a+pr.
		\]
		Выразим $c\cdot1$, $c\cdot a$, $c\cdot a^2$ через $1,a,a^2$:
		\begin{align*}
			c\cdot1&=1+a+a^2,\\
			c\cdot a&=a^3+a^2+a=-r+(1-q)\,a+(1-p)\,a^2,\\
			c\cdot a^2&=a^4+a^3+a^2=(pr-r)+(pq-q-r)\,a+(p^2-p-q+1)\,a^2.
		\end{align*}
		В последней строке мы сложили
		$a^4=pr+(pq-r)a+(p^2-q)a^2$, $a^3=-r-qa-pa^2$ и $a^2$.

		Запишем это матрично. Пусть $w=(1,\,a,\,a^2)^{T}$ и
		\[
			A=\begin{pmatrix}
				1 & 1 & 1\\
				-r & 1-q & 1-p\\
				pr-r & pq-q-r & p^2-p-q+1
			\end{pmatrix}.
		\]
		Тогда три полученных равенства означают, что $Aw=c\,w$, то есть
		\[
			(A-cE)\,w=0.
		\]
		Вектор $w\ne0$ (его первая координата равна $1$), поэтому однородная
		система $(A-cE)w=0$ имеет ненулевое решение, и
		\[
			\det(cE-A)=0.
		\]
		Значит, $c$~--- корень характеристического многочлена
		$\chi(x)=\det(xE-A)$. Для матрицы $3\times3$
		\[
			\chi(x)=x^3-\operatorname{tr}A\cdot x^2+S\,x-\det A,
		\]
		где $S$~--- сумма главных миноров второго порядка. Вычислим эти
		величины.

		\emph{След.}
		\[
			\operatorname{tr}A=1+(1-q)+(p^2-p-q+1)=p^2-p-2q+3.
		\]

		\emph{Главные миноры.}
		\begin{align*}
			M_{12}&=1\cdot(1-q)-1\cdot(-r)=1-q+r,\\
			M_{13}&=1\cdot(p^2-p-q+1)-1\cdot(pr-r)=p^2-p-q+1-pr+r,\\
			M_{23}&=(1-q)(p^2-p-q+1)-(1-p)(pq-q-r)\\
			&=p^2-p-q+1-pq+q^2+r-pr.
		\end{align*}
		Складывая,
		\[
			S=M_{12}+M_{13}+M_{23}
			=2p^2-2p-pq+q^2-2pr-3q+3r+3.
		\]

		\emph{Определитель.} Раскладываем по первой строке:
		\begin{align*}
			\det A&=M_{23}
				-\bigl((-r)(p^2-p-q+1)-(1-p)(pr-r)\bigr)
				+\bigl((-r)(pq-q-r)-(1-q)(pr-r)\bigr)\\
			&=M_{23}+(pr-qr)+(r^2-pr+r)\\
			&=p^2+q^2+r^2-pq-pr-qr-p-q+2r+1.
		\end{align*}

		Итак, $c$~--- корень многочлена
		\[
			\boxed{\begin{aligned}
				h(x)={}&x^3-(p^2-p-2q+3)\,x^2\\
				&+(2p^2-2p-pq+q^2-2pr-3q+3r+3)\,x\\
				&-(p^2+q^2+r^2-pq-pr-qr-p-q+2r+1).
			\end{aligned}}
		\]
	\end{proof}

	\noindent
	\emph{Проверка.} Пусть $f=x^3-2$ ($p=q=0$, $r=-2$), $a=\sqrt[3]{2}$.
	Формулы дают $h(x)=x^3-3x^2-3x-1$. Непосредственно, используя $a^3=2$,
	\[
		c^2=a^4+2a^3+3a^2+2a+1=3a^2+4a+5,\qquad
		c^3=c\cdot c^2=12a^2+15a+19,
	\]
	откуда
	\[
		c^3-3c^2-3c-1=(12-9-3)a^2+(15-12-3)a+(19-15-3-1)=0.
	\]

	%==================================================
	\section*{Задача 5}

	\noindent
	\textbf{Условие.} Дан неприводимый многочлен $f\ne x$ над $\mathbb{Q}$,
	причём $x$ и $x^2$~--- два из его корней. Докажите, что $x$~--- корень
	из единицы.

	\medskip
	\noindent
	Чтобы не путать корень с переменной, обозначим его через $\alpha$.
	Итак, $f(\alpha)=f(\alpha^2)=0$; докажем, что $\alpha^N=1$ при некотором
	натуральном $N$. Корни рассматриваем в $\mathbb{C}$. Условие $f\ne x$
	понимается с точностью до ненулевого множителя.

	\begin{proof}
		\emph{Шаг 1: все корни $f$ отличны от нуля.} Если $f(0)=0$, то
		$x\mid f$, и из неприводимости $f$ следует $f=cx$, что исключено.
		Значит, $f(0)\ne0$; в частности, $\alpha\ne0$.

		\emph{Шаг 2: $f\mid f(x^2)$ в $\mathbb{Q}[x]$.} Положим
		$g(x)=f(x^2)\in\mathbb{Q}[x]$. Тогда $g(\alpha)=f(\alpha^2)=0$.
		Наибольший общий делитель $d=\gcd(f,g)$ является делителем
		неприводимого $f$, поэтому $d=1$ или $d=f$ (с точностью до
		константы). Если $d=1$, то найдутся $u,v\in\mathbb{Q}[x]$ с
		$uf+vg=1$; подставляя $x=\alpha$, получаем $0=1$~--- противоречие.
		Значит, $d=f$, то есть $f\mid g$:
		\[
			f(x^2)=f(x)\,h(x),\qquad h\in\mathbb{Q}[x].
		\]

		\emph{Шаг 3: квадрат корня~--- снова корень.} Пусть $R$~--- множество
		корней $f$ (оно конечно: $|R|\leqslant\deg f$). Если $\beta\in R$, то
		\[
			f(\beta^2)=f(\beta)\,h(\beta)=0,
		\]
		то есть $\beta^2\in R$. Поскольку $\alpha\in R$, по индукции
		\[
			\alpha,\ \alpha^2,\ \alpha^4,\ \alpha^8,\ \ldots,\ \alpha^{2^k},\ \ldots\in R.
		\]

		\emph{Шаг 4: вывод.} Множество $R$ конечно, а в последовательности
		$\alpha^{2^k}$ бесконечно много членов, поэтому найдутся $i<j$ с
		\[
			\alpha^{2^i}=\alpha^{2^j}.
		\]
		Так как $\alpha\ne0$, можно поделить на $\alpha^{2^i}$:
		\[
			\alpha^{2^j-2^i}=1,
		\]
		причём $N=2^j-2^i>0$. Следовательно, $\alpha$~--- корень из единицы.
	\end{proof}

\end{document}